\chapter{Measuring Variability and Metastability}
\label{ch:measuring-variability-metastability}

\textit{A theory of clustered spiking networks is only useful if its variables reproduce measured signatures: spike-count variability, irregularity, slow autocorrelations, cluster dwell times, transition paths, dimensionality, correlations, and spectra. This chapter defines those measurements and the main traps in estimating them.}

% ============================================================
\section{Spike-count variance}
\label{sec:spike-count-variance}
% ============================================================

Let $N_i^{(q)}(T)$ be the number of spikes from neuron or population $i$ in trial $q$ during a counting window of duration $T$. The sample mean and variance across $R$ trials are
%
\begin{equation}
  \bar N_i(T)=\frac{1}{R}\sum_{q=1}^R N_i^{(q)}(T),
  \qquad
  s_i^2(T)=\frac{1}{R-1}\sum_{q=1}^R
  \left(N_i^{(q)}(T)-\bar N_i(T)\right)^2.
  \label{eq:sample-count-variance}
\end{equation}
%
The mean measures response magnitude; the variance measures trial-to-trial reliability at that window size.

For spontaneous simulations with no trial structure, one often uses non-overlapping windows as pseudo-trials. That is valid only if the process is approximately stationary over the analyzed period. Slow drift can masquerade as high count variance.

% ============================================================
\section{Fano factor}
\label{sec:fano-factor-measurement}
% ============================================================

The Fano factor normalizes count variance by count mean:
%
\begin{equation}
  F_i(T)=\frac{\Var[N_i(T)]}{\E[N_i(T)]}.
  \label{eq:fano-measurement}
\end{equation}
%
For a Poisson process, $F_i(T)=1$. Values below one indicate counts more regular than Poisson; values above one indicate overdispersion from bursting, slow rate fluctuations, shared latent state, or mixtures of response modes.

The window duration $T$ is part of the definition. Short windows emphasize spike-generation irregularity. Long windows include slow cluster-state fluctuations. In clustered networks, $F(T)$ increasing with $T$ is often a signature that slow mesoscopic dynamics are contributing to variability.

% ============================================================
\section{CV and local CV}
\label{sec:cv-local-cv}
% ============================================================

For interspike intervals $T_m=t^{m+1}-t^m$, the coefficient of variation is
%
\begin{equation}
  \mathrm{CV} =
  \frac{\sqrt{\Var[T_m]}}{\E[T_m]}.
  \label{eq:isi-cv-measurement}
\end{equation}
%
Near-periodic firing has $\mathrm{CV}<1$, Poisson firing has $\mathrm{CV}=1$, and bursty or heavy-tailed firing has $\mathrm{CV}>1$.

The local CV reduces sensitivity to slow rate changes by comparing adjacent intervals:
%
\begin{equation}
  \mathrm{CV}_{\mathrm{local}}
  =
  \frac{3}{M-1}\sum_{m=1}^{M-1}
  \left(
  \frac{T_{m+1}-T_m}{T_{m+1}+T_m}
  \right)^2.
  \label{eq:local-cv}
\end{equation}
%
The numerator measures local interval contrast; the denominator normalizes by local firing scale. This statistic helps separate fast spike-time irregularity from slow gain modulation.

In E/I clustered attractor models, an important possibility is high spike irregularity within each state but reduced slow count variability after stimulus onset. CV and Fano factor therefore answer different questions.

% ============================================================
\section{Mean-matched Fano factor}
\label{sec:mean-matched-fano}
% ============================================================

Fano factors can change simply because mean rates change. Mean matching controls this confound by comparing conditions using only time bins, neurons, or trials with matched mean-count distributions.

A practical procedure is:
\begin{enumerate}[leftmargin=*]
  \item estimate mean count for every neuron/time bin/condition unit;
  \item choose a common mean-count range shared by the compared conditions;
  \item subsample or weight units so the mean-count histograms match;
  \item compute Fano factors on the matched set.
\end{enumerate}

Mean matching is essential for variability quenching. If stimulus onset increases firing rates and the Fano factor drops, the drop should not be interpreted as reduced latent variability until rate-related sampling effects are controlled.

% ============================================================
\section{Autocorrelation functions}
\label{sec:autocorrelation-functions}
% ============================================================

For a stationary scalar signal $x(t)$, the normalized autocorrelation is
%
\begin{equation}
  \rho_x(\tau)=
  \frac{\E[(x(t)-\mu_x)(x(t+\tau)-\mu_x)]}{\Var[x(t)]}.
  \label{eq:autocorrelation}
\end{equation}
%
It measures how much the present predicts the future after lag $\tau$. The signal might be a smoothed population rate, a cluster activation indicator, or a low-dimensional latent variable.

Fast irregular spiking produces short autocorrelation. Metastable cluster switching produces slow tails because the current state remains predictive for the dwell time of that state. A useful summary is the integrated autocorrelation time,
%
\begin{equation}
  \tau_{\mathrm{int}}=\int_0^\infty \rho_x(\tau)\,\dd\tau,
  \label{eq:integrated-autocorrelation-time}
\end{equation}
%
provided the integral is well estimated and the process is stationary.

% ============================================================
\section{Cluster lifetime distributions}
\label{sec:cluster-lifetime-distributions}
% ============================================================

Define a cluster as active when its rate or activity exceeds a threshold for at least a minimal duration. The lifetime $L$ is the contiguous time spent in that active state.

The survival function
%
\begin{equation}
  S_L(t)=\mathbb{P}\{L>t\}
  \label{eq:lifetime-survival}
\end{equation}
%
is often more stable to plot than a histogram. Exponential survival suggests approximately memoryless escape from a state. Heavy tails suggest heterogeneous barriers, slow adaptation, latent inputs, or a mixture of state types.

Lifetime estimates depend on smoothing, thresholds, and merging rules for brief interruptions. Report those choices. A cluster-state claim is weak if it disappears under small changes in detection parameters.

% ============================================================
\section{Transition-triggered averages}
\label{sec:transition-triggered-averages}
% ============================================================

A transition-triggered average aligns trajectories at detected switch times. If $t_m$ is the time of transition $m$, then
%
\begin{equation}
  \bar x_{\mathrm{TTA}}(\tau)
  =
  \frac{1}{M}\sum_{m=1}^M x(t_m+\tau).
  \label{eq:transition-triggered-average}
\end{equation}
%
The average can be computed for source cluster rate, target cluster rate, inhibitory activity, adaptation, synaptic current, or a low-dimensional order parameter.

The goal is to see the route through state space. Does inhibition drop before a cluster turns on? Does the target cluster ramp gradually or jump abruptly? Do transitions pass through a common saddle-like region? These questions connect measured switching to the dynamical-systems picture in \cref{ch:dynamical-systems-essentials}.

Avoid circular interpretation: if transitions are detected from $x(t)$ itself, the aligned average of $x$ will necessarily show a crossing. More informative variables are those not used for detection.

% ============================================================
\section{Dimensionality of cluster activity}
\label{sec:dimensionality-cluster-activity}
% ============================================================

Let $\Sigma$ be the covariance matrix of cluster activities across time or trials. If $\lambda_1,\ldots,\lambda_K$ are its eigenvalues, the participation-ratio dimensionality is
%
\begin{equation}
  d_{\mathrm{PR}}
  =
  \frac{\left(\sum_{k=1}^K \lambda_k\right)^2}
       {\sum_{k=1}^K \lambda_k^2}.
  \label{eq:participation-ratio}
\end{equation}
%
If one mode dominates, $d_{\mathrm{PR}}\approx 1$. If many modes contribute equally, $d_{\mathrm{PR}}$ is large.

In clustered networks, dimensionality measures how many independent cluster combinations are explored. A stimulus that biases the network into fewer states should reduce dimensionality even if within-state spiking remains irregular.

The covariance should be computed on variables that match the theoretical claim: binary cluster activity indicators for attractor occupancy, smoothed rates for population dynamics, or residuals after subtracting stimulus-locked means for trial-to-trial variability.

% ============================================================
\section{Pairwise correlations}
\label{sec:pairwise-correlations}
% ============================================================

For spike counts $N_i$ and $N_j$ in the same window, the noise correlation is
%
\begin{equation}
  c_{ij}=
  \frac{\operatorname{Cov}(N_i,N_j)}
       {\sqrt{\Var[N_i]\Var[N_j]}}.
  \label{eq:pairwise-count-correlation}
\end{equation}
%
Correlations should usually be computed after subtracting condition means, otherwise stimulus tuning can create signal correlations that are not noise correlations.

Clustered networks predict structure in correlations: neurons in the same active cluster should share slow fluctuations, while neurons in competing clusters may be anticorrelated when only a few clusters can be active at once. Inhibitory clustering can change this pattern by stabilizing local E/I balance and reducing shared runaway excursions.

Small pairwise correlations can still have large population effects when many neuron pairs share the same sign. Always ask whether correlations are independent pair facts or signatures of a low-dimensional shared state.

% ============================================================
\section{Power spectra}
\label{sec:power-spectra}
% ============================================================

The power spectrum of a zero-mean signal $x(t)$ is the Fourier transform of its autocovariance:
%
\begin{equation}
  S_x(f)=\int_{-\infty}^{\infty}
  C_x(\tau)e^{-2\pi i f\tau}\,\dd\tau.
  \label{eq:power-spectrum-definition}
\end{equation}
%
Slow switching produces excess low-frequency power. Oscillations produce peaks at nonzero frequencies. White noise gives a flat spectrum over the frequency range where the approximation holds.

For simulations, spectra are useful because they separate mechanisms by time scale. A clustered attractor model should show power at dwell-time frequencies; an E/I oscillatory instability should show a sharper peak; independent spike noise should dominate high frequencies.

Use consistent preprocessing: bin width, smoothing, detrending, tapering, and total recording duration all affect the estimate.

% ============================================================
\section{Practical pitfalls in estimating these quantities}
\label{sec:estimation-pitfalls}
% ============================================================

The common pitfalls are statistical rather than philosophical:
\begin{enumerate}[leftmargin=*]
  \item \textbf{Nonstationarity}: slow drift inflates variance, autocorrelation, and low-frequency power.
  \item \textbf{Window choice}: Fano factors and correlations change with bin size.
  \item \textbf{Rate confounds}: variability measures must often be mean-matched.
  \item \textbf{Threshold dependence}: cluster lifetimes depend on smoothing and active-state criteria.
  \item \textbf{Pooling heterogeneous neurons}: mixing cells with different rates or tuning can create artificial overdispersion.
  \item \textbf{Too few transitions}: lifetime tails and transition-triggered averages need enough detected switches.
  \item \textbf{Autocorrelated samples}: overlapping windows are not independent trials.
  \item \textbf{Detection circularity}: do not validate a transition using only the variable that defined it.
  \item \textbf{Finite simulation length}: rare escapes and low-frequency spectra require long runs.
  \item \textbf{Unreported preprocessing}: smoothing, filtering, and detrending must be stated.
\end{enumerate}

A good measurement section in a rotation artifact should specify the observable, window, conditioning, estimator, uncertainty estimate, and robustness checks.

\begin{chaptersummary}

\begin{center}
\renewcommand{\arraystretch}{1.5}
\begin{tabularx}{\textwidth}{@{}l X X@{}}
  \toprule
  \textbf{Measurement} & \textbf{Formula} & \textbf{What it diagnoses} \\
  \midrule
  Count variance &
  $\Var[N(T)]$ &
  Trial-to-trial reliability in a chosen window \\
  Fano factor &
  $F=\Var[N]/\E[N]$ &
  Overdispersion or quenching beyond mean rate \\
  ISI CV &
  $\sqrt{\Var[T]}/\E[T]$ &
  Spike-time irregularity \\
  Local CV &
  $\frac{3}{M-1}\sum((T_{m+1}-T_m)/(T_{m+1}+T_m))^2$ &
  Local irregularity less sensitive to slow rate drift \\
  Autocorrelation &
  $\rho_x(\tau)$ &
  Persistence of rates, states, or latent variables \\
  Lifetime survival &
  $S_L(t)=\mathbb{P}\{L>t\}$ &
  Dwell-time statistics of active clusters \\
  Dimensionality &
  $d_{\mathrm{PR}}=(\sum\lambda_k)^2/\sum\lambda_k^2$ &
  Number of activity modes explored \\
  Correlation &
  $c_{ij}=\operatorname{Cov}(N_i,N_j)/(\sigma_i\sigma_j)$ &
  Shared variability within or across clusters \\
  Spectrum &
  $S_x(f)=\mathcal{F}\{C_x(\tau)\}$ &
  Time scales: slow switching, oscillations, or fast noise \\
  \bottomrule
\end{tabularx}
\end{center}

\end{chaptersummary}

\begin{conceptquestions}
  \item Why is the counting-window duration part of the definition of a Fano factor?
  \item How can CV remain high while stimulus onset reduces Fano factor?
  \item What does mean matching control for in variability-quenching analyses?
  \item Why are cluster lifetime distributions more informative than just the mean lifetime?
  \item What variable would you use for a transition-triggered average if the transition was detected from cluster rate?
  \item How would you distinguish low-frequency power from slow nonstationary drift?
\end{conceptquestions}

\clearpage
